% !PRINTED 136
\weberpage{136}
%<<CONT
is then contained in none of the groups $Q_1, Q_2, \dots, Q_{q-1}$, because it was assumed that $Q$ has no divisor in common with any of these groups whose degree is greater than $p^{\gamma}$. In the same way we can find a divisor $R_1$ of $Q_1$ of degree $p^{\gamma+1}$, of which $Q_0$ is a normal divisor, and which is contained in none of the groups $Q, Q_2, \dots, Q_{q-1}$. Now $R$ as well as $R_1$ is contained in $P$. We consider the smallest group $S$ that contains $R$ and $R_1$ at the same time, which is in any case contained in $P$ (the least common multiple of $R$ and $R_1$). The degree of this group $S$ cannot be a power of $p$, since otherwise, by Theorem II., \S.\,30, $S$, and hence $R$ and $R_1$, would have to be contained in one of the groups $Q, Q_1, \dots, Q_{q-1}$. But $R$ is contained only in $Q$, $R_1$ only in $Q_1$; hence this assumption is inadmissible. The degree of $S$ must therefore be of the form $p^{\lambda} q$.

But $\lambda$ cannot be smaller than $\alpha$; for, since a group of degree $p^{\gamma+1}$ is contained in $S$, namely $R$, we have

\weq{4}{\lambda \geqq \gamma+1 > \beta.}

By the hypothesis made in $\alpha)$, however, no divisor of degree $p^{\lambda} q$ is contained in $P$ in which $\lambda$ is at the same time greater than $\beta$ and smaller than $\alpha$. By (4), therefore, necessarily $\lambda = \alpha$, i.\,e. $S$ is identical with $P$.

Now $Q_0$ is a normal divisor of $R$ and $R_1$. If, then, we collect all elements $c$ of $P$ that satisfy the condition

\[ c^{-1}\,Q_0\,c = Q_0 \]

into a group, this group contains $R$ as well as $R_1$, and is therefore the whole group $P$. Hence $Q_0$ is also a normal divisor of $P$. Since by (3) $\gamma$ is greater than zero, the degree of $Q_0$ is greater than 1, and we come upon a contradiction with the assumption that $P$ is simple.

\begin{center}
\S.\,34.

\gesperrt{Simple Groups.}
\end{center}

The theorems that we have come to know in the preceding sections admit fairly far-reaching conclusions about the nature of groups. They show that, at least for the lower degrees, the metacyclic (and cyclic) groups decidedly predominate. The
%CONT>>
